\origpage{39}{GRAPH MINORS. XIV}
\textit{Let $P_1, \dots, P_\beta$ be disjoint paths of $H$ from $V(C_2)$ to $V(C_1)$, such that for $1 \le i \le \beta$, $P_i$ has initial vertex $v_i \in V(C_2)$ and $U(P_i) \cap (\Delta_2 - bd(\Delta_2)) = \emptyset$. Then $\{v_1, \dots, v_\beta\}$ is free with respect to the tangle of order $\theta - 4\lambda - 2$ obtained from $\mathscr{T}$ by clearing the $\lambda$-zone $\Delta_2 - bd(\Delta_2)$.}

\textit{Proof.} Let $H' = H \cap (\Sigma - (\Delta_2 - bd(\Delta_2)))$, and let $\mathscr{T}'$ be the tangle of order $\theta - 4\lambda - 2$ in $H'$ obtained by clearing $\Delta_2 - bd(\Delta_2)$. Suppose that $\{v_1, \dots, v_\beta\}$ is not free with respect to $\mathscr{T}'$, and choose $(A, B) \in \mathscr{T}'$ of order $< \beta$ with $v_1, \dots, v_\beta \subseteq V(A)$. This is possible, for certainly $\beta \le \theta - 4\lambda - 2$. Since $|V(A \cap B)| < \beta$ and $P_1, \dots, P_\beta$ are mutually disjoint, there exists $j$ with $1 \le j \le \beta$ such that $V(P_j) \cap V(A \cap B) = \emptyset$. Since $P_j$ has an end in $V(A)$ and $P_j$ is a connected subgraph of $H'$, it follows that $V(P_j) \subseteq V(A) - V(B)$.

Since $|V(A \cap B)| < \beta$ there exists $i$ with $1 \le i \le \beta/2$ such that $|V(C_i \cap A \cap B)| \le 1$, because $C_1, \dots, C_{\beta/2}$ are mutually disjoint. Since $U(C_i)$ bounds a disc $\Delta$ in $\Delta_1$ including $\Delta_2$, and $P_j$ is a path from $V(C_2)$ to $V(C_1)$, it follows that $V(C_i \cap P_j) \ne \emptyset$. Since $V(P_j) \subseteq V(A) - V(B)$ we deduce that $V(C_i) \not\subseteq V(B)$. Since $|V(C_i \cap A \cap B)| \le 1$ it follows that $C_i \subseteq A$ and $|V(C_i \cap B)| \le 1$. By (7.4) there is a circuit $D$ of a radial drawing $K'$ of $H'$ such that $V(D) \cap V(H') \subseteq V(A \cap B)$ and $U(C_i) \subseteq ins'(D)$, where $ins'$ is the inside function of $\mathscr{T}'$. Now $\Delta - bd(\Delta)$ is a $(\theta - \beta - 4\lambda - 2)$-zone around $z$ (in $H, \mathscr{T}$), and hence is a $(\theta - \beta - 4\lambda - 2)$-zone around the region $\Delta_2 - bd(\Delta_2)$ (in $H', \mathscr{T}'$). Since $|E(D)| < 2\beta$ and $bd(\Delta) \subseteq ins'(D)$, it follows from (7.6) applied to $H', \mathscr{T}'$ that $\Delta \subseteq ins'(D)$. In particular, $U(D) \cap \Delta_2 \subseteq bd(\Delta_2)$. Choose a radial drawing $K$ of $H$ so that $K' \setminus v$ is a subgraph of $K$, where $v$ is the unique vertex of $K'$ in $\Delta_2 - bd(\Delta_2)$; then $D \subseteq K$. Since $|U(D) \cap V(H)| < \beta \le \theta - 4\lambda - 2$, it follows that $ins'(D) = ins(D)$. Since $z \in \Delta_2 \subseteq \Delta \subseteq ins'(D) = ins(D)$, we deduce from (iv) that $ins(D) \subseteq \Delta_1 - bd(\Delta_1)$. But there is a vertex of $P_j$ in $ins(D)$, and $U(P_j) \cap U(D) = \emptyset$ because $V(P_j) \cap V(A \cap B) = \emptyset$, and so $U(P_j) \subseteq ins(D) \subseteq \Delta_1 - bd(\Delta_1)$. Yet $P_j$ has an end in $C_1$, a contradiction.

Hence $\{v_1, \dots, v_\beta\}$ is free with respect to $\mathscr{T}'$, as required. $\square$

\sectionhead{8. \textsc{Segregation} \textsc{In} A \textsc{Rural} \textsc{Neighbourhood}}

Now we develop some lemmas about societies. If $\Omega$ is a cyclic permutation, let us write $\bar{\Omega} = \{v_1, \dots, v_n\}$, numbered in order under $\Omega$. An $\Omega$-interval is a set $\{v_j : i \le j \le k\}$ where $0 \le i \le k \le n$, or the complement in $\bar{\Omega}$ of such a set. (This definition is independent of which element of $\bar{\Omega}$ is chosen to be $v_1$). A \textit{transaction} in a society $(G, \Omega)$ is a set $\mathscr{P}$ of mutually disjoint paths of $G$, each with distinct ends both in $\bar{\Omega}$, such that there is an $\Omega$-interval containing exactly one end of each member of $\mathscr{P}$. If $\mathscr{P}$ is a transaction, $\bigcup \mathscr{P}$ denotes the union of the members of $\mathscr{P}$, and $V(\mathscr{P})$ denotes the set of ends
